\chapter{Reading Molecular Models}

\begin{kaichang}
Open a medicine leaflet, a popular-science article, or a later chapter of this book, and you may encounter letters joined by short lines; zigzags with a letter or two at their ends; colored balls connected by little sticks; or bulging ``bubbles'' crowded together. All claim to depict the same thing: a molecule.

Why draw one molecule in four or five completely different ways? Which is its true appearance?

None of them, and all of them. This appendix teaches five common molecular portraits: Lewis structures, skeletal formulas, wedge-and-dash formulas, ball-and-stick models, and space-filling models. Each drawing makes a trade: it preserves some information while deliberately discarding other information. Reading them is less about memorizing drawing rules than understanding the bargain. What does it want you to see? What does it hide to keep you from being distracted?
\end{kaichang}

\section{A Drawing Is Not the Molecule}

All five portraits share a principle: \textbf{they are human conventions for representation}. Real atoms are not colored; black, red, and blue balls distinguish elements. Chemical bonds are neither sticks nor lines but arrangements of electrons between atoms (Chapters 17 and 19). Molecules are not static photographs either: they constantly vibrate, rotate, and bend. Every model resembles a map, flattening a three-dimensional city and omitting alleyways to make the main roads clear at a glance. The first question is what the drawing trades for what.

The following five models run from least ink to most space. Each moves closer to the molecular outline while making the interior harder to see. We will meet them in that order.

\section{Lewis Structures: The Electron Ledger}

Lewis structures, whose history and limitations Chapter 19 discussed, are the most accountant-like models. They answer one question: \textbf{where are the valence electrons?}

There are three reading rules:

\begin{itemize}[itemsep=2pt]
\item Element symbols represent nuclei plus inner-shell electrons.
\item One line between atoms represents a shared electron pair, a covalent bond; two lines are a double bond, three a triple bond.
\item Dots beside atoms show nonbonding lone pairs, often crucial in reactions and the basis for Chapter 32's Lewis acid--base definition.
\end{itemize}

For water, put oxygen between two hydrogens with one line to each, and a pair of dots above and below oxygen. In two seconds you read two \chem{O-H} single bonds and two lone pairs. For carbon dioxide, put carbon between two oxygens with two lines on each side and two pairs of dots on each oxygen. The carbon--oxygen bonds are double, and each atom's electron account balances. This is the model's strength: it records every valence electron clearly enough to audit the ledger.

Line count itself carries information. In the oxygen you breathe, \chem{O_2}, the two oxygens share two lines, a double bond. In the air's most abundant gas, \chem{N_2}, the nitrogens share three, a triple bond. This locks them together so tightly that separation takes great effort, explaining why nitrogen is approximately $78\%$ of air yet hardly reacts, and why plants need rhizobia or fertilizer factories to use it (Chapter 86). A tiny Lewis structure already writes down why nitrogen is so unreactive.

\textbf{What it preserves}: connectivity, bond multiplicity, lone-pair counts, and which atom carries charge. These allow many predictions about reaction mechanisms.

\textbf{What it sacrifices}: almost all geometry. A Lewis drawing may make water straight or right-angled, but its actual \chem{O-H} angle is approximately $104.5^{\circ}$. Chapter 20 explained that this bend determines polarity and, in turn, nearly all water's unusual properties (Chapter 23). Dots also suggest stationary beads, whereas electrons are diffuse clouds, as Chapter 19 emphasized. For benzene, whose electrons spread around a ring, no single Lewis structure suffices; several resonance structures must take turns (Chapters 19 and 41).

A common test is drawing charged ions. Enclose the whole structure in square brackets and place its charge at the upper right. Hydroxide, \chem{OH^-}, water's relative, has one oxygen--hydrogen line and three oxygen lone pairs, enclosed in brackets with a minus sign. The ion has one extra electron. Carbonate, \chem{CO_3^{2-}}, needs more than one drawing: the double bond can take turns among three oxygens, motivating resonance structures (Chapter 19) and exposing the ledger's limits.

Use it when counting electrons, checking charges, or writing mechanisms. It is an account book, not a photograph.

\section{Skeletal Formulas: Organic Shorthand}

Organic textbooks and new-drug news are full of zigzag lines. These skeletal, or bond-line, formulas appeared with carbon chains in Chapter 36. Hydrocarbons contain too many carbons and hydrogens to write every one, so chemists invented this shorthand.

Again there are three rules, and almost every beginner misreads them:

\begin{itemize}[itemsep=2pt]
\item Every \textbf{line end and corner} represents a carbon atom, whose symbol is omitted.
\item Hydrogens bonded to carbon are also omitted. Assume carbon always makes four bonds (Chapter 36); any bond not occupied by another atom goes to hydrogen.
\item Atoms other than carbon and hydrogen, such as oxygen, nitrogen, and chlorine, must be labeled explicitly. Hydrogens attached to them are usually written too, as in $-\chem{OH}$.
\end{itemize}

Ethanol, for example, is a short two-segment bend ending in $\chem{OH}$. Two segments, one corner, and three letters describe all of \chem{C_2H_6O}'s connectivity: two carbons carrying five hydrogens, with a hydroxyl group attached to the terminal carbon. For acetic acid from Chapter 40, vinegar's sour component, add a carbon--oxygen double bond beside ethanol's shorthand, and the behavioral difference becomes visually clear.

Two further details often appear in questions. First, rings: a closed hexagon represents a six-membered ring. Add a circle or three alternating short lines and it is a benzene ring (Chapter 41), whose six carbons share delocalized electrons. Do not ask which particular edge is the double bond; the circle says the electrons are spread out. Second, parallel lines indicate double bonds, such as the carbon--oxygen double bond in carboxyl groups. Try aspirin's skeletal formula: a benzene ring with adjacent substituents, one carrying a carboxyl group and the other a small ester group. With these two tools, you can read the connectivity in many medicine-leaflet diagrams.

\textbf{What it preserves}: carbon-skeleton connectivity, whether straight, branched, or cyclic, and functional-group positions, Chapter 37's reaction interfaces. Zigzags do more than save ink: carbon chains rotate around single bonds (Chapter 36), and a zigzag roughly represents a common extended posture.

\textbf{What it sacrifices}: carbon--hydrogen detail, requiring mental restoration of hydrogens; actual atomic sizes; exact bond angles, since a drawn $120^{\circ}$ corner is not a real approximately $109.5^{\circ}$ angle; and all three-dimensional information. Skeletal drawings flatten molecules onto paper, so readers must remember their spatial nature. Ordinary skeletal formulas lose chirality entirely, requiring the next model.

Use it to follow changes in large carbon skeletons, recognize functional groups, or quickly compare connectivity, as in Chapter 36's isomers.

\section{Wedges and Dashes: Out of the Page}

Skeletal formulas pretend molecules are flat; wedges and dashes expose the illusion by adding two special lines:

\begin{itemize}[itemsep=2pt]
\item A \textbf{solid triangular wedge} shows a bond projecting out of the page toward you.
\item A \textbf{hashed wedge}, short parallel strokes increasing in length, shows a bond receding behind the page.
\item An ordinary thin line shows a bond in the plane of the page.
\end{itemize}

Their most familiar use is Chapter 42's chirality. A carbon bonded to four different groups has two spatial arrangements that are mirror images but cannot be superimposed, like your hands. Ordinary flat formulas make them look identical. Draw one group on a solid wedge, one on a hashed wedge, and two on plain lines, and left and right become unmistakable. For lactic acid, one arrangement occurs in muscles aching after exercise, and the other is its mirror image. Your enzymes recognize only one because they are themselves chiral pockets (Chapters 42 and 49).

A practical reading method is to find the focal atom where the solid and hashed wedges meet, usually the chiral center. Identify which groups point toward you, away, and within the page. Some textbooks also widen wedges from a narrow end to a broad end, the broad end nearer you. The page is only a reference frame: rotate the whole drawing mentally while retaining the relative spatial relationships. This trains Chapter 42's three-dimensional vision.

\textbf{What it preserves}: front-to-back group orientations at key atoms, or stereochemical information. In medicines this can be life-critical: Chapter 42 showed that mirror-image molecules can have very different effects, one therapeutic and the other ineffective or harmful.

\textbf{What it sacrifices}: only the focal region gains spatial relationships; the rest remains flattened. Wedges show relative directions, not exact angles or distances. A subtler trap is constant single-bond rotation (Chapter 36): the drawing freezes one of countless conformations, not the molecule's entire life.

Use it for chirality, comparing mirror images, and tracking changes in group orientation during reactions.

\section{Ball-and-Stick Models: Visible Frameworks}

Now we move from paper symbols to physical likenesses. The most famous model uses colored balls for atoms and sticks for covalent bonds. Predrilled holes guide sticks into the correct angles, such as carbon's approximately $109.5^{\circ}$ tetrahedral angle. Such models appear in almost every chemistry-laboratory photograph.

Colors are conventional, with manufacturer variations: carbon black, hydrogen white, oxygen red, nitrogen blue, sulfur yellow, chlorine green. Memorize this: \textbf{these are conventions, not atomic colors}. Atoms are hundreds of times smaller than visible wavelengths, so color has no meaning for them.

\begin{qiaoliang}
Why do methane's four bonds make approximately $109.5^{\circ}$ angles? Not aesthetics, but energy accounting. Four bonding electron pairs surround carbon, repel because they are negative, and lower total energy by separating. In a plane, four directions can only form a $90^{\circ}$ cross. Point them out of the page toward a regular tetrahedron's vertices and every pair makes approximately $109.5^{\circ}$, the roomiest arrangement. Chapter 20 used this electron-pair-repulsion idea to explain bent water and pyramidal ammonia. Molecular blueprints are written in electron-repulsion accounts.
\end{qiaoliang}

\textbf{What it preserves}: connectivity, bond angles, and overall geometry. Turn the model in your hand and Chapter 20's principle becomes tangible: molecular shape determines behavior.

\textbf{What it sacrifices}: first, scale. Long rigid sticks suggest open gaps between atoms, whereas bonded electron clouds closely overlap; real molecules contain neither sticks nor the apparent space they occupy. Chapter 17 dismantled that picture. Second, electrons disappear: no lone pairs or diffuse clouds remain visible.

Use it to introduce shapes and angles, or ask whether a molecule can fit a particular pocket.

\section{Space-Filling Models: Molecular Bulk}

Ball-and-stick models shrink atoms and lengthen bonds; space-filling models do the reverse. Expand each atom to its actual relative size, press the spheres together, and remove the sticks. The resulting bumpy bubbles most closely resemble molecular outlines among these five models. Folded proteins in Chapter 48, enzyme pockets in Chapter 49, and DNA grooves in Chapter 66 often use them.

An order-of-magnitude estimate explains their greater honesty. Typical covalent bonds are around $10^{-10}$ meters long; a \chem{C-H} bond is approximately \ud{1.1\times 10^{-10}}{m}. Carbon's radius is itself more than half this order of magnitude. Bonded atoms' center-to-center distance is therefore smaller than their combined radii: the spheres must overlap deeply, leaving no visible gap. Two small balls staring across a stick do not exist at real scale; overlapping space-filling bubbles better approximate occupied space.

\textbf{What it preserves}: occupied volume and the surface's shape and contours. Whether a drug slips into an enzyme pocket, a substrate passes through a membrane channel (Chapter 51), or water contacts a protein surface are contact questions only a space-filling model can answer.

\textbf{What it sacrifices}: bonds and buried atoms disappear. A space-filling protein looks like cauliflower, revealing no connectivity; angles, bond orders, and lone pairs are gone. It is honest about outline and silent about what lies beneath.

Use it when molecular contact matters: recognition, binding, blockage, or passage.

\section{One Molecule, Five Readings}

Definitions alone are not enough. Follow a real molecule through all five: alanine, one of twenty amino acids in your proteins (Chapters 48 and 69). Its file lists a central carbon bonded to an amino group, $-\chem{NH_2}$; a carboxyl group, $-\chem{COOH}$; a methyl group, $-\chem{CH_3}$; and hydrogen.

\textbf{First, Lewis structure}. Draw every atom and count electron pairs: four single bonds at the central carbon, one nitrogen lone pair, one carbon--oxygen double bond in the carboxyl group, and two lone pairs on each oxygen. The books balance, and nitrogen's lone pair explains its ability to accept a proton, part of why amino acids act as both acids and bases (Chapter 32).

\textbf{Second, skeletal formula}. Omit carbon symbols and carbon-bound hydrogens. A short bent line joins methyl to the central carbon, from which two strokes lead to $\chem{NH_2}$ and $\chem{COOH}$. Three seconds reveal the skeleton and three functional-group positions.

\textbf{Third, wedges and dashes}. Four different groups make the central carbon chiral. Put amino on a solid wedge toward you, hydrogen on a hashed wedge behind, and methyl and carboxyl on ordinary lines. That is one alanine; swapping wedge directions gives its mirror image. Your proteins use almost exclusively one form (Chapter 42).

\textbf{Fourth, ball-and-stick}. Black carbon, white hydrogen, blue nitrogen, and red oxygen balls connect through sticks at approximately $109.5^{\circ}$. Rotate the model to see amino and carboxyl spatial relationships with an intuition flat drawings cannot provide.

\textbf{Fifth, space-filling}. Expand all atoms to real sizes and let them overlap. Sticks vanish, leaving bumpy bubbles with a little blue projection at the amino group. Their shape determines whether alanine fits a particular enzyme pocket (Chapter 49).

Five readings answer five questions. Whenever alanine appears in a book, you will see one of these portraits. Recognize which, and you know what the author wants you to notice.

\section{Comparing the Five Models}

\begin{table}[htbp]
\centering
\begin{tabular}{llll}
\toprule
Model & Question & Preserves & Sacrifices \\
\midrule
Lewis structure & \shortstack[l]{Where are\\valence electrons?} & \shortstack[l]{Connectivity, bond\\orders, lone pairs} & \shortstack[l]{3D shape,\\electron diffusion} \\
Skeletal formula & \shortstack[l]{How is carbon\\connected?} & \shortstack[l]{Carbon connectivity,\\functional groups} & \shortstack[l]{C and H, real angles,\\three dimensions} \\
\shortstack[l]{Wedge-and-\\dash} & \shortstack[l]{Which points\\forward or back?} & \shortstack[l]{Key spatial\\orientations} & \shortstack[l]{Exact angles, other\\rotating\\conformations} \\
Ball-and-stick & \shortstack[l]{What is the\\framework geometry?} & \shortstack[l]{Angles, overall\\configuration} & \shortstack[l]{Real scale,\\electron clouds} \\
Space-filling & \shortstack[l]{How much space\\does it occupy?} & \shortstack[l]{Real size,\\surface shape} & \shortstack[l]{Bonds, interior\\atoms, connectivity} \\
\bottomrule
\end{tabular}
\caption{Each molecular model trades information. Colors, sizes, sticks, and lines are conventions, not atoms' actual appearance.}
\end{table}

The table's key conclusion is simple: \textbf{there is no best model, only a model suited to the present question}. Chapter 38 uses skeletal formulas for organic reactions, Chapter 42 wedges for chirality, and Chapter 49 space-filling models for enzyme pockets. The book has always changed portraits to suit its questions. In any scientific article, first identify the portrait, then ask whether it suits the question.

\begin{shiyan}
\textbf{Feel molecular shape: gummy-candy models}

Materials: gummies in several colors, or small modeling-clay balls, and a box of toothpicks.

Procedure: 1. Build water from one oxygen ball and two hydrogen balls, joining them with toothpicks bent into an obtuse angle of approximately $105^{\circ}$; do not place the hydrogens opposite each other. 2. Build carbon dioxide from one carbon and two oxygens, this time arranging them in a straight line. 3. Build methane from one carbon and four hydrogens, trying to make all pairwise toothpick angles as equal as possible. They naturally spread into a three-dimensional frame that will not flatten into a cross.

What to observe: why will methane's four bonds not remain in one plane? Hint: see the earlier bridge box. Put water beside carbon dioxide and predict which has positive and negative ends and which does not. Chapters 20 and 22 reveal the answer.

Safety: toothpicks are sharp; do not point them at yourself or anyone else. Used gummies collect dust and bacteria: \textbf{do not eat them}.
\end{shiyan}

\begin{wuqu}
``Which molecular model looks most real? I want to learn the truest one.'' The premise is wrong. Space-filling models best resemble the outline but tell nothing about bonds or interiors; Lewis structures resemble molecules least yet reveal electron details. Asking which is truer is like comparing a world map and a subway map. One preserves shape and distance, the other transfer connections; bring a world map into a subway and you still get lost. A model's truth is not photographic resemblance but \textbf{whether its retained information answers your question}. Likewise, a blue ball does not mean a blue atom. Atoms have no color; blue is the illustrator's uniform for nitrogen.
\end{wuqu}

\begin{zhengju}
Molecular models were once mocked predictions, not merely classroom charts. In 1874, the young Dutch chemist van 't Hoff was unknown when he published a booklet proposing that carbon's four bonds point to a regular tetrahedron's vertices. The evidence was indirect. If four bonds lay in one plane, dichloromethane, \chem{CH_2Cl_2}, should have two differently arranged forms, but chemists had found only one. A tetrahedron explained that uniqueness. The prominent chemist Kolbe publicly mocked it as speculation on paper.

Yet the paper model kept making falsifiable predictions: carbon carrying four different groups should have left- and right-handed forms. Over subsequent decades, optical activity in lactic acid and other molecules, the rotation of polarized light discussed in Chapter 42, fulfilled the prediction. Nearly forty years later, in 1912, Laue and the Braggs pioneered X-ray crystallography, letting people first directly ``see'' actual atomic arrangements in crystals. Angles and distances matched the paper model. Van 't Hoff therefore received the first Nobel Prize in Chemistry in 1901. The memorable point is that a good molecular model is not a picture but a prediction machine: it dares to describe unseen things and then submits to evidence.
\end{zhengju}

\section{Shared Boundaries}

Finally, consider where all five fail. In delocalized molecules such as benzene (Chapter 41), one-stick-one-bond pictures mislead: those six electron pairs belong to no single pair of atoms. Metals (Chapter 21) have no distinct molecules to draw. Proteins with thousands of atoms turn even space-filling models into unrecognizable cauliflower, forcing structural biologists to invent sixth and seventh representations, such as ribbons (Chapter 48).

Whenever you see a molecular drawing, ask what it wants you to see and what it hides to prevent distraction. That habit is what it means to understand molecular models, and every diagram this book presents next.

\section*{Exercises}

\begin{sikaoti}
\item Glucose is \chem{C_6H_{12}O_6}. Compare its Lewis and space-filling representations. Which lets you count hydroxyl groups, $-\chem{OH}$, and which helps judge whether it fits a membrane channel? What does this comparison demonstrate?
\item An unlabeled corner appears in a skeletal formula. Which atom is it, and how many hydrogens does it carry? Which rule lets you answer?
\item Can a standard skeletal formula distinguish lactic acid from its mirror image? Which model is needed, and why can the distinction be life-critical in medicine? Recall Chapter 42.
\item Long thin sticks exaggerate atomic separation. If a carbon ball has diameter 2 centimeters and its bonded hydrogen ball approximately 1.5 centimeters, how far apart should their centers be at real scale? Hint: the balls must overlap deeply. Which of the five models gets this proportion right?
\item Find a popular-science article about a medicine or nutrient. Identify its molecular-model type, then explain what the author wants you to see, what is hidden, and which other model would restore that information.
\end{sikaoti}
